\section{Theta--zeta branch: modular inversion, Mellin scale, completion, and primes}

\subsection{Jacobi theta modular inversion}

Define
\begin{equation}
\vartheta(t)=\sum_{n\in\mathbb Z}e^{-\pi n^2t},\qquad t>0.
\end{equation}
The modular inversion is
\begin{equation}
\boxed{\vartheta(t)=t^{-1/2}\vartheta(1/t).}
\label{eq:theta-modular}
\end{equation}
Status: \status{STANDARD\_EXACT} \cite{DLMFTheta}. Introduce \(u=\log t\). Then multiplicative inversion becomes additive reflection \(u\mapsto-u\). A symmetrised profile can therefore be chosen even in log-scale.

\subsection{Mellin transform to zeta}

A standard theta Mellin relation yields the gamma--zeta factor; one convenient normalisation is
\begin{equation}
\boxed{
\int_0^\infty[\vartheta(t)-1]t^{s/2-1}\,dt
=2\pi^{-s/2}\Gamma(s/2)\zeta(s),
\qquad \RePart(s)>1.
}
\label{eq:theta-mellin}
\end{equation}
Equivalent theta-function Mellin formulae are tabulated by DLMF \cite{DLMFThetaMellin}.

For \(t=e^u\), Mellin kernels take the form
\begin{equation}
 t^{-s}=e^{-su}=e^{-\sigma u}e^{-i\tau u},
\qquad s=\sigma+i\tau.
\end{equation}
Thus the imaginary part of \(s\) is a Fourier-like phase frequency conjugate to logarithmic scale.

\subsection{Completed zeta and the half-axis}

Define
\begin{equation}
\xi(s)=\frac12s(s-1)\pi^{-s/2}\Gamma(s/2)\zeta(s).
\end{equation}
Then
\begin{equation}
\boxed{\xi(s)=\xi(1-s).}
\label{eq:xi-reflection}
\end{equation}
Status: \status{STANDARD\_EXACT} \cite{DLMFZetaReflection}. Because \(\xi(\bar s)=\overline{\xi(s)}\), the anti-linear involution
\begin{equation}
J(s)=1-\bar s
\end{equation}
has fixed set
\begin{equation}
\boxed{J(s)=s\iff \RePart(s)=\frac12.}
\end{equation}
This fixed-set statement does not prove that all non-trivial zeros lie on that set.

\subsection{Euler product and the prime alphabet}

For \(\RePart(s)>1\),
\begin{equation}
\boxed{\zeta(s)=\prod_{p\in\mathbb P}(1-p^{-s})^{-1}.}
\label{eq:euler-product}
\end{equation}
Status: \status{STANDARD\_EXACT} \cite{DLMFZetaEuler}. Every positive integer has a unique prime-valuation vector
\begin{equation}
\mathbf v(n)=\{v_p(n)\}_{p\in\mathbb P},
\qquad
n=\prod_p p^{v_p(n)},
\end{equation}
and therefore
\begin{equation}
\boxed{\log n=\sum_p v_p(n)\log p.}
\label{eq:prime-log}
\end{equation}
Only finitely many terms are non-zero for any fixed integer.

This identifies a common multiplicative coordinate system. Zeta packages independent prime powers through the Euler product in its convergence half-plane; Collatz, developed in the next branch, generates a deterministic trajectory through the same integer valuation lattice.

\subsection{Critical-line prime characters}

For a single prime and \(s=\tfrac12+i\tau\),
\begin{equation}
\boxed{p^{-s}=p^{-1/2}e^{-i\tau\log p}.}
\label{eq:prime-critical-character}
\end{equation}
Thus a finite prime character has amplitude \(p^{-1/2}\) and phase \(-\tau\log p\). This identity is valid for an individual finite factor for any complex \(s\).

\begin{claim}{Convergence boundary -- \status{MANDATORY}.}
Equation~\eqref{eq:euler-product} is a directly convergent infinite product only for \(\RePart(s)>1\). The project must not numerically multiply an unregularised infinite Euler product on \(\RePart(s)=1/2\) and call the result \(\zeta(s)\). Critical-line work must use the analytically continued \(\zeta/\xi\) or explicitly finite, receipt-bound prime-mode truncations.
\end{claim}

\INTERP
The prime representation supplies the missing multiplicative ``alphabet'' between zeta and Collatz. It closes a representation bridge: both systems can be expressed using \(\log p\) and \(v_p\). It does not yet close a spectral identity between the Riemann operator sought by Hilbert--P\'olya-type programmes and a Collatz transfer/Koopman operator.
